\documentclass[11pt,openany]{book}
\usepackage[margin=1in]{geometry}
\usepackage[T1]{fontenc}
\usepackage{lmodern}
\usepackage{microtype}
\usepackage{amsmath,amssymb,amsthm}
\usepackage{booktabs}
\usepackage{enumitem}
\usepackage{hyperref}
\usepackage{titlesec}

\hypersetup{colorlinks=true,linkcolor=black,citecolor=black,urlcolor=blue}

\titleformat{\chapter}[display]
  {\normalfont\huge\bfseries}
  {\chaptertitlename\ \thechapter}{16pt}{\Huge}
\titlespacing*{\chapter}{0pt}{30pt}{20pt}

\titleformat{\section}
  {\normalfont\Large\bfseries}{\thesection}{1em}{}

\titleformat{\subsection}
  {\normalfont\large\bfseries}{\thesubsection}{1em}{}

\setlength{\parskip}{0.4em}
\setlength{\parindent}{0pt}

\title{LLM Systems Reader: Math Refresher}
\author{Companion to \emph{Large Language Models as Engineered Systems}}
\date{\today}

\begin{document}

\frontmatter

\maketitle
\tableofcontents

\chapter*{How to Use This Companion}
\addcontentsline{toc}{chapter}{How to Use This Companion}

Use this as a lookup tool. Read only the chapter you need for the reader chapter you are working through. If a chapter points to an earlier refresher chapter, follow that pointer instead of re-reading repeated material. The chapter numbers here match those of the main reader exactly.

\mainmatter

\chapter*{Preface: Probability Building Blocks}
\addcontentsline{toc}{chapter}{Preface: Probability Building Blocks}

\section*{How to read this preface}
The chapters ahead assume undergraduate probability fluency that may have faded. This preface does not replace a textbook; it reinstalls a \emph{Strang-style chain}: each subsection states what problem the next definition solves, names where the idea entered mathematics or engineering, and ends with a \emph{mental model} you will reuse when reading LLM systems notation.

Probability here means Kolmogorov-style measure on a sample space. Random variables compress events into summaries (numbers or vectors). Expectation averages those summaries under a distribution. Information quantities reinterpret probabilities as coding costs. Large-sample intuition connects finite training batches to population quantities. Later refresher chapters pointer-reference these steps instead of repeating them.

\section*{Step 1 --- Sample space, events, and axioms}
\textbf{Problem.} We must assign consistent numbers to uncertain outcomes so algebraic manipulation respects logic.

\textbf{Construction (compact).} A sample space $\Omega$ collects mutually exclusive elementary outcomes. An event is a subset $A\subseteq\Omega$ (in the $\sigma$-algebra $\mathcal{F}$ when rigor demands). A probability measure $P$ obeys $P(\Omega)=1$, $P(A)\ge 0$, and countable additivity on disjoint unions.

\textbf{Insight.} Uncertainty becomes bookkeeping: combine events with set operations; probabilities inherit monotone laws that mirror everyday implication.

\textbf{Lineage.} Kolmogorov's axioms (1933) unified earlier paradox-prone intuitions into measure theory---still the lingua franca of ML papers.

\textbf{Mental model.} Before manipulating symbols, write down what $\Omega$ \emph{is}. Ambiguous $\Omega$ produces ambiguous conditioning later---especially when ``token'' vs ``character'' vs ``byte'' outcomes differ.

\section*{Step 2 --- Conditional probability}
\textbf{Problem.} Observations arrive partially: we learn $B$ happened and must revise beliefs about $A$.

\textbf{Definition.} For $P(B)>0$,
\[
  P(A\mid B)=\frac{P(A\cap B)}{P(B)}.
\]
Equivalently $P(A\cap B)=P(A\mid B)P(B)$.

\textbf{Insight.} Conditioning \emph{restricts and renormalises} the world to the slice where $B$ is true.

\textbf{Lineage.} Formalised within Kolmogorov's programme; conceptual ancestors include Bernoulli/Huygens gaming mathematics and later Bayesian narratives.

\textbf{Mental model.} Language-model ``next token given context'' \emph{is} conditional probability: $B$ is ``we saw prefix $x_{<t}$,'' $A$ is ``next symbol equals $x_t$.''

\section*{Step 3 --- Chain rule for probabilities}
\textbf{Problem.} High-dimensional joints look unmanagable unless we factor them along an ordering.

\textbf{Identity.} Always $P(A,B)=P(A\mid B)P(B)$. Iterating along an ordered sequence gives
\[
  P(x_{1:T})=\prod_{t=1}^{T} P(x_t\mid x_{<t}),
\]
provided each conditional is defined (positive-probability prefixes).

\textbf{Insight.} Sequential prediction needs no extra axiom beyond Step~2---factorisation is algebraic honesty.

\textbf{Lineage.} Elementary consequence of definitions; Markov (1913) showcased sequential dependence in text before computers existed.

\textbf{Mental model.} Each factor answers ``what fraction of surviving probability mass lands on this symbol \emph{after} seeing the left context?'' Misordered dependencies confuse causality with notation.

\section*{Step 4 --- Bayes' theorem (compact)}
\textbf{Problem.} Sometimes we observe effects and must infer latent causes or revise hypotheses.

\textbf{Form.} With $P(E)>0$, $P(H\mid E)=\dfrac{P(E\mid H)\,P(H)}{P(E)}$.

\textbf{Insight.} Invert conditioning direction by paying the evidence normaliser $P(E)$.

\textbf{Lineage.} Bayes; Laplace championed inferential uses.

\textbf{Mental model.} Pure LM training rarely expands Bayes in full, but calibration papers and Bayesian interpretations of weight decay/noise often invoke posterior language---recognise when authors silently swap $\theta$ and data roles.

\section*{Step 5 --- Random variables and expectation}
\textbf{Problem.} Events are verbose; we want numerical summaries whose averages behave linearly.

\textbf{Discrete expectation.}
\[
  \mathbb{E}[f(X)]=\sum_{x\in\mathcal{X}} p(x)\,f(x).
\]
Linearity: $\mathbb{E}[a f(X)+b g(X)] = a\,\mathbb{E}[f(X)] + b\,\mathbb{E}[g(X)]$ (no independence needed).

\textbf{Insight.} Expectation is the ``center of mass'' of outcomes weighted by probability---the quantity laws of large numbers later approximate.

\textbf{Lineage.} Nineteenth-century mechanics analogies $\to$ measure-theoretic integral (Lebesgue).

\textbf{Mental model.} Empirical cross-entropy is a finite-sample proxy for $\mathbb{E}_{x\sim\text{data}}[-\log q_\theta(x)]$; stochastic gradients perturb that proxy.

\section*{Step 6 --- Variance and sums (sketch)}
\textbf{Problem.} Expectation hides spread; optimisation noise depends on second moments.

\textbf{Definitions.} $\mathrm{Var}(X)=\mathbb{E}[(X-\mathbb{E}[X])^2]$; for independent $X,Y$, $\mathrm{Var}(X+Y)=\mathrm{Var}(X)+\mathrm{Var}(Y)$.

\textbf{Insight.} Independent minibatch gradients sum variances---batch scaling trades compute against noise amplitude.

\textbf{Lineage.} Classical; Chebyshev's inequality links variance to tail bounds.

\textbf{Mental model.} When inputs correlate inside a batch, $\mathrm{Var}(\sum)$ gains covariance terms---clip-heavy regimes acknowledge heavy tails violating naive Gaussian pictures.

\section*{Step 7 --- Law of large numbers (operational)}
\textbf{Problem.} Finite datasets nonetheless approximate unseen distributions---when is that morally justified?

\textbf{Qualitative statement.} Under IID or weakened dependence, sample averages $\frac{1}{n}\sum_{i=1}^n f(X_i)$ converge to $\mathbb{E}[f(X)]$ as $n\to\infty$ (modes differ: weak vs strong LLN).

\textbf{Insight.} Training losses reference an underlying population quantity even though only finite corpora arrive.

\textbf{Lineage.} Early probabilists (Bernoulli); modern proofs measure-theoretic.

\textbf{Mental model.} Engineering monitors finite-$n$ deviation (bias, variance, drift); LLM corpora violate IID---LLN intuition guides, assumptions rarely hold verbatim.

\section*{Step 8 --- Maximum likelihood}
\textbf{Problem.} Choose parameters $\theta$ so observed data look plausible under $P_\theta$.

\textbf{Objective.} $\hat\theta\in \arg\max_\theta \prod_{i=1}^n P_\theta(\text{datum}_i)$; equivalently maximise $\sum_i \log P_\theta(\text{datum}_i)$.

\textbf{Insight.} ML estimation aligns probabilities with frequencies inside the model class---not proof of truth outside that class.

\textbf{Lineage.} Fisher sharpened theory in the 1920s--1930s; modern deep learning uses regularised variants.

\textbf{Mental model.} Products become sums via $\log$---derivatives and curvature numerics prefer sums; ``negative log-likelihood'' is the minimisation convention.

\section*{Step 9 --- Shannon entropy}
\textbf{Problem.} Quantify uncertainty with a single scalar aligned with coding limits.

\textbf{Definition (discrete).} $H(p)=-\sum_x p(x)\log p(x)$ (nats if $\ln$, bits if $\log_2$).

\textbf{Insight.} Optimal lossless codes for draws from $p$ require $\approx H(p)$ symbols per draw on average---entropy is not mysticism, it is asymptotic compression cost.

\textbf{Lineage.} Shannon's \emph{A Mathematical Theory of Communication} (1948).

\textbf{Mental model.} Never compare entropies computed in different log bases without conversion; tokenisation changes $|\mathcal{X}|$ and hence typical $H$ magnitudes.

\section*{Step 10 --- Cross-entropy and Kullback--Leibler divergence}
\textbf{Problem.} Models $q$ disagree with reference $p$; quantify mismatch in bits/nats.

\textbf{Definitions.}
\begin{align*}
  H(p,q) &= -\sum_x p(x)\log q(x),\\
  D_{\mathrm{KL}}(p\|q) &= \sum_x p(x)\log\frac{p(x)}{q(x)}.
\end{align*}
\textbf{Decomposition.} $H(p,q)=H(p)+D_{\mathrm{KL}}(p\|q)$.

\textbf{Insight.} Cross-entropy splits unavoidable uncertainty ($H(p)$) plus excess coding penalty from using $q$ when truth is $p$ (the KL term).

\textbf{Lineage.} Shannon coding theorems motivate $H(p,q)$; Kullback and Leibler (1951) formalised $D_{\mathrm{KL}}$.

\textbf{Mental model.} Minimising sample cross-entropy pushes $q$ toward empirical token frequencies \emph{inside the softmax parameterisation}; KL penalties in alignment explicitly drag policies toward references---same symbolic animal, different rhetorical frame.

\section*{Where Chapter 1 picks up}
Chapter~1 of this refresher specialises Steps~2--3 to autoregressive tokens, Steps~8--10 to softmax logits and calibration diagnostics, and binds Step~7 qualitatively to dataset-scale claims. Re-read this preface whenever conditioning or information identities feel like unmotivated algebra.


\section{Chapter 1 Refresher: Statistical Learning Foundations for Large Language Models --- Probability, Likelihood, and Calibration}

\subsection*{Setting the Scene}
Chapter~1 of the reader uses a specific branch of probability theory: the branch that turns uncertain events into measurable predictive statements. In practice, the reader will use four ingredients repeatedly: conditional probability, the chain rule, expectation, and logarithmic scoring.

The path to this toolkit was gradual and each step added a new capability:
\begin{itemize}[leftmargin=1.5em]
  \item \textbf{Axiomatic probability (Kolmogorov):} gave a rigorous language for uncertainty over events.
  \item \textbf{Conditional probability and chain decompositions:} made sequential prediction mathematically legal.
  \item \textbf{Information theory (Shannon):} interpreted uncertainty as expected coding burden.
  \item \textbf{Relative entropy (Kullback--Leibler):} measured model mismatch relative to the true distribution.
  \item \textbf{Statistical learning formulation:} converted these quantities into trainable optimization objectives.
\end{itemize}
This sequence is exactly the reasoning path used in modern language modeling.

The critical insight for Chapter~1 is that language modeling does not require a new probability theory. It requires careful use of known results in a specific sequence. We use conditional probability to represent next-token prediction, chain rule to represent whole sequences, logarithms to make optimization stable and additive, and information-theoretic quantities to interpret what loss values mean.

The purpose of this refresher is to rebuild that sequence gently. The target is not symbolic performance. The target is that each mathematical move has a clear intent in the reasoning chain.

\subsection*{What You Need To Recall Precisely}
\begin{itemize}[leftmargin=1.5em]
  \item \textbf{Conditional probability as prediction primitive.} If $A$ and $B$ are events with $P(B)>0$, then
  \[
    P(A\mid B)=\frac{P(A\cap B)}{P(B)}.
  \]
  For language modeling, $A$ is ``next token equals $x_t$'' and $B$ is ``context is $x_{<t}$''.

  \item \textbf{Chain rule as sequence constructor.}
  \[
    p(x_{1:T}) = \prod_{t=1}^{T} p(x_t \mid x_{<t}).
  \]
  This is not a modeling trick; it is repeated application of conditional probability.

  \item \textbf{Likelihood to optimization objective.} Maximum likelihood chooses parameters that maximize probability of observed data. Because products are awkward numerically, we apply log:
  \[
    \arg\max_\theta \prod_i a_i(\theta)=\arg\max_\theta \sum_i \log a_i(\theta).
  \]
  Then we minimize negative average log-likelihood (NLL).

  \item \textbf{Expectation as average under a distribution.} If $X$ is discrete:
  \[
    \mathbb{E}[f(X)] = \sum_x p(x)f(x).
  \]
  This lets us interpret empirical losses as estimates of population quantities.

  \item \textbf{Entropy, cross-entropy, and KL in one chain.}
  For distributions $p$ (reference/true) and $q$ (model):
  \[
    H(p) = -\sum_x p(x)\log p(x),\quad
    H(p,q) = -\sum_x p(x)\log q(x),
  \]
  \[
    D_{\mathrm{KL}}(p\|q)=\sum_x p(x)\log\frac{p(x)}{q(x)}.
  \]
  Then
  \[
    H(p,q) = H(p) + D_{\mathrm{KL}}(p\|q).
  \]
  This identity explains why minimizing cross-entropy is equivalent to minimizing KL divergence up to an additive term independent of the model.

  \item \textbf{Softmax as probability map from logits:}
  \[
    \mathrm{softmax}(z)_i = \frac{e^{z_i}}{\sum_j e^{z_j}}.
  \]
  \textbf{Shift invariance (numerical stability):}
  \[
    \mathrm{softmax}(z) = \mathrm{softmax}(z - c\mathbf{1}) \quad \text{for any scalar } c.
  \]
\nobreak
  Usually $c=\max_i z_i$ to avoid overflow.

  \item \textbf{Perplexity as exponentiated cross-entropy.}
  \[
    \mathrm{PPL}=\exp(H)
  \]
  when $H$ is in nats. Perplexity is interpretable only with fixed tokenization and log base.

  \item \textbf{Calibration distinction.} A model can assign high probability and still be wrong frequently. Loss and calibration answer different questions.
\end{itemize}

\subsection*{How These Results Build on Each Other}
\begin{enumerate}[leftmargin=1.5em]
  \item Conditional probability defines next-token prediction.
  \item Chain rule extends next-token predictions to full-sequence probabilities.
  \item Likelihood turns full-sequence probabilities into a training objective.
  \item Log transform converts objective from product form to additive form.
  \item Empirical averaging converts sample objective to dataset objective.
  \item Cross-entropy/KL decomposition gives interpretation of what minimizing the objective does.
  \item Softmax provides a valid distribution over vocabulary for each context.
  \item Calibration analysis audits whether confidence values are trustworthy after fitting.
\end{enumerate}

\subsection*{Reminder Glossary (Term by Term)}
\begin{itemize}[leftmargin=1.5em]
  \item \textbf{Random variable:} a map from outcomes to numerical values.
  \item \textbf{Conditional probability:} probability under restricted information.
  \item \textbf{Likelihood:} probability of observed data viewed as function of parameters.
  \item \textbf{Log-likelihood:} logarithm of likelihood; additive over samples/tokens.
  \item \textbf{NLL:} negative log-likelihood; minimized in training.
  \item \textbf{Entropy $H(p)$:} uncertainty of true distribution.
  \item \textbf{Cross-entropy $H(p,q)$:} expected code length when coding $p$ with model $q$.
  \item \textbf{KL divergence $D_{\mathrm{KL}}(p\|q)$:} mismatch penalty from using $q$ when truth is $p$.
  \item \textbf{Logits:} pre-softmax scores in $\mathbb{R}^{|V|}$.
  \item \textbf{Softmax:} normalization from logits to a probability simplex.
  \item \textbf{Perplexity:} exponentiated average NLL/cross-entropy (with fixed log base).
  \item \textbf{Calibration:} correspondence between predicted confidence and observed accuracy.
  \item \textbf{ECE:} binned estimate of calibration gap.
\end{itemize}

\subsection*{Worked Example (Non-Toy, End-to-End Reasoning)}
Assume a validation slice with $N=1000$ predicted tokens. We observe:
\begin{itemize}[leftmargin=1.5em]
  \item Mean NLL $=1.10$ nats.
  \item Mean confidence of top-1 predictions $=0.78$.
  \item Actual top-1 accuracy $=0.64$.
\end{itemize}
Now reason in the chain used in Chapter~1:

\textbf{Step 1: Loss to perplexity.}
\[
\mathrm{PPL}=\exp(1.10)\approx 3.00.
\]
Interpretation: average uncertainty corresponds to about three equiprobable choices per token on this slice.

\textbf{Step 2: Loss does not equal calibration.}  
Despite moderate loss, confidence-accuracy gap is
\[
\Delta_{\mathrm{cal}} = 0.78 - 0.64 = 0.14.
\]
The model is overconfident by 14 percentage points on average.

\textbf{Step 3: Why this matters for downstream chapters.}  
Chapter 9 (alignment) can reduce harmful outputs while still leaving calibration issues.  
Chapter 10 (RAG/evaluation) can improve factual grounding while confidence remains miscalibrated.  
Therefore the distinction between fit quality and confidence quality is operational, not philosophical.

\textbf{Step 4: Practical intervention path.}
\begin{enumerate}[leftmargin=1.5em]
  \item Keep objective fit diagnostics (NLL/perplexity).
  \item Add calibration diagnostics (reliability bins, ECE).
  \item Apply temperature scaling on held-out data.
  \item Re-check both loss and calibration after scaling.
\end{enumerate}

\subsection*{Intuition Checkpoints}
\begin{itemize}[leftmargin=1.5em]
  \item A lower NLL does \emph{not} guarantee calibrated confidence.
  \item KL is a divergence, not a metric; direction matters.
  \item Perplexity is only comparable under matched tokenization and log base.
  \item ECE is a binned approximation and depends on binning and sample size.
\end{itemize}

\subsection*{If You Get Stuck}
\begin{itemize}[leftmargin=1.5em]
  \item Rewrite every probability expression in words before manipulating symbols.
  \item Track whether each quantity is theoretical, empirical estimate, or model output.
  \item Check if the reasoning step is algebraic (identity) or statistical (estimation).
\end{itemize}

\subsection*{References and What To Use Them For}
\begin{itemize}[leftmargin=1.5em]
  \item \textbf{Murphy, \emph{Probabilistic Machine Learning}} --- conditional probability, expectation, estimation viewpoint.
  \item \textbf{Bishop, \emph{Pattern Recognition and Machine Learning}} --- likelihood-based learning, cross-entropy objectives.
  \item \textbf{Cover--Thomas, \emph{Elements of Information Theory}} --- entropy, cross-entropy, KL interpretation.
  \item \textbf{Guo et al. (2017)} --- practical calibration diagnostics and temperature scaling.
  \item \textbf{Wasserman, \emph{All of Statistics}} --- fast refresh on estimation, uncertainty, and empirical risk language.
\end{itemize}

\chapter{Chapter 2 Refresher: Optimization, Generalization, and Scaling Intuition --- Optimization and Scaling}

\section*{Setting the Scene}
Chapter~2 consumes the gradient of an empirical risk built in Chapter~1 (Preface Steps~5--8): we optimise an expectation estimated from batches. The novelty is size: billions of parameters force \emph{cheap, noisy} directional probes rather than exact curvature solves.

\paragraph{Step A --- Local linear model.}
\textbf{Problem.} Full loss landscapes are unwieldy; near $\theta_t$ we approximate $\mathcal{L}(\theta_t+\Delta)$ by first-order Taylor expansion.
\textbf{Insight.} Gradient descent assumes meaningful linear prediction within chosen step lengths.
\textbf{Lineage.} Cauchy (1847) gradient methods; modern proofs invoke convex-like intuitions rarely satisfied literally in deep nets.
\textbf{Mental model.} If curvature tightens, shrink steps or add momentum filters---otherwise linearisation lies.

\paragraph{Step B --- Stochastic gradients.}
\textbf{Problem.} Exact $\nabla\mathcal{L}$ over full corpora is unaffordable.
\textbf{Insight.} Subsample to inject noise whose variance trades memory for stability (Preface Step~6).
\textbf{Lineage.} Robbins--Monro (1951); Bottou's large-scale ML synthesis.
\textbf{Mental model.} Treat minibatch gradient as true gradient + zero-mean perturbation; diagnose spikes as breakdown of bounded-variance stories.

\paragraph{Step C --- Adaptive signal processing.}
\textbf{Problem.} Parameter blocks receive gradients at wildly different scales (embeddings vs logits vs biases).
\textbf{Insight.} Running moment estimates rescale updates elementwise while preserving descent intent at smooth regions.
\textbf{Lineage.} Polyak momentum; Duchi/Hazan/Singer (Adagrad); Kingma/Ba (Adam); Loshchilov/Hutter (AdamW).
\textbf{Mental model.} Adaptive methods behave like learned preconditioners---trust them until drift/outliers falsify stationarity.

\paragraph{Step D --- Numerical hygiene.}
\textbf{Problem.} Low-precision tensors amplify underflow/overflow.
\textbf{Insight.} Loss scaling preserves gradient magnitude while exploiting hardware throughput.
\textbf{Lineage.} NVIDIA mixed-precision training recipes (2017 era onward).
\textbf{Mental model.} Clip gradients when tail events violate FP16 comfort zones.

\paragraph{Step E --- Macroscopic empirical laws.}
\textbf{Problem.} Engineers must budget parameters \emph{and} tokens jointly.
\textbf{Insight.} Iso-loss contours summarize decades of runs---Preface Step~7 interpreted informally at massive $n$.
\textbf{Lineage.} Kaplan et al.\ (2020); Hoffmann et al.\ / Chinchilla (2022).
\textbf{Mental model.} Scaling fits extrapolate poorly across dataset contamination regimes---treat coefficients as local Taylor surrogates.

\paragraph{Connective thread.}
Exact gradients are expensive; minibatches trade bias for variance: write $\hat g_t=\nabla_\theta \mathcal{L}_t(\theta_t)+\varepsilon_t$ with $\mathbb{E}[\varepsilon_t]=0$. Under i.i.d.\ sampling, $\mathrm{Var}(\varepsilon_t)$ scales roughly like $1/B$ with batch size $B$, which is why larger batches smooth updates but cost memory. Momentum and adaptive moments are filters on that noise; clipping guards rare spikes that violate the ``bounded variance'' story.

\section*{Notation Introduced in This Chapter}
\begin{itemize}[leftmargin=1.5em]
  \item $\theta$: model parameters.
  \item $\mathcal{L}(\theta)$: objective (loss) as function of parameters.
  \item $\nabla\mathcal{L}(\theta)$: gradient vector.
  \item $\eta_t$: learning rate at step $t$.
  \item $\hat g_t$: stochastic (minibatch) gradient estimate.
  \item $m_t, v_t$: first and second moment EMAs.
  \item $\beta_1,\beta_2$: EMA decay coefficients.
  \item $\tau$: clipping threshold.
  \item $N,D$: model/data scale variables used in scaling-law fits.
\end{itemize}

\section*{What You Need To Recall Precisely}
\begin{itemize}[leftmargin=1.5em]
  \item \textbf{Gradient as first-order local direction.} If $\mathcal{L}(\theta)$ is differentiable, then for small $\Delta$:
  \[
    \mathcal{L}(\theta+\Delta)\approx \mathcal{L}(\theta) + \nabla \mathcal{L}(\theta)^\top \Delta.
  \]
  Choosing $\Delta=-\eta\nabla\mathcal{L}$ gives first-order descent.

  \item \textbf{Stochastic gradient descent (SGD):}
  \[
    \theta_{t+1}=\theta_t-\eta_t \nabla \mathcal{L}(\theta_t).
  \]
  In practice, $\nabla \mathcal{L}$ is replaced by minibatch estimate $\hat g_t$.

  \item \textbf{Momentum as exponential averaging.}  
  With $m_t=\beta m_{t-1} + (1-\beta)\hat g_t$, updates use smoothed direction $m_t$.
  Think of $m_t$ as a low-pass filter on the gradient signal: high-frequency minibatch noise is attenuated while consistent downhill curvature survives, similar in spirit to a heavy ball continuing through shallow oscillations in the loss surface.

  \item \textbf{Bias correction for startup EMA.}  
  Because $m_0=0$, early EMAs are biased toward zero---the running average ``forgets'' that gradients were nonzero before time zero. Adam uses debiased estimates
  \[
    \hat m_t=\frac{m_t}{1-\beta_1^t},\qquad \hat v_t=\frac{v_t}{1-\beta_2^t}.
  \]
  so that if $\hat g$ were constant, $\hat m_t$ would recover that constant after a few steps instead of ramping up spuriously slowly.

  \item \textbf{AdamW intuition (brief).}
  Classic Adam entangles weight decay with adaptive preconditioning in ways that distort the intended $\ell_2$ penalty. AdamW applies weight decay \emph{directly} to $\theta$ (decoupled from the adaptive gradient transform), while still using $\hat m_t,\hat v_t$ to precondition the loss gradient step.


  \item \textbf{Gradient clipping rule:}
  \[
    g \leftarrow g \cdot \min\left(1,\frac{\tau}{\|g\|}\right).
  \]
  This bounds update magnitude without changing direction when clipping is active.

  \item \textbf{Mixed precision with loss scaling.}  
  Forward/backward often run in bf16/fp16 for throughput; tiny gradients can flush to subnormals or zero in those formats. Multiply the scalar loss by a large factor $s$ before backward pass so chain-rule gradients inherit scale; divide accumulated gradients by $s$ before the optimizer applies updates (often fp32)---recovering usable magnitude without changing the mathematically correct descent direction.

  \item \textbf{Scaling-law interpretation.}  
  Empirical forms such as
  \[
    L(N,D)=L_\infty + aN^{-\alpha} + bD^{-\beta}
  \]
  summarise measured iso-loss contours across runs when architecture and data pipelines stay comparable. Use them as \emph{budget planners}: given fixed compute $C\approx 6ND$ tokens-parameter conventions from Kaplan/Chinchilla-style analyses, pick $(N,D)$ pairs predicted to sit near the elbow rather than extrapolate blindly across regimes where data quality or objectives shifted.
\end{itemize}

\section*{How These Results Build on Each Other}
\begin{enumerate}[leftmargin=1.5em]
  \item First-order approximation motivates gradient descent.
  \item Minibatch estimation introduces noise, requiring variance management.
  \item Momentum/adaptive moments stabilize noisy trajectories.
  \item Bias correction fixes startup distortion in adaptive moments.
  \item Clipping and precision controls handle numerical and tail-event instability.
  \item Scaling fits summarize many such runs for planning decisions.
\end{enumerate}

\section*{Reminder Glossary (Term by Term)}
\begin{itemize}[leftmargin=1.5em]
  \item \textbf{Gradient:} local rate-of-change vector for loss.
  \item \textbf{Learning rate:} step size in parameter space.
  \item \textbf{Minibatch noise:} estimation error from finite sample gradients.
  \item \textbf{Momentum/EMA:} smoothed running estimate of gradient direction.
  \item \textbf{Second moment estimate:} running estimate of squared gradients.
  \item \textbf{Bias correction:} startup debiasing of EMAs.
  \item \textbf{Gradient clipping:} norm-bound safeguard against spikes.
  \item \textbf{Warmup:} gradual LR increase to stabilize early dynamics.
  \item \textbf{Scaling law:} empirical relation among model size, data, compute, and loss.
  \item \textbf{AdamW:} Adam-style adaptive steps with weight decay applied directly to parameters (decoupled from gradient preconditioning).
\end{itemize}

\section*{Worked Example (Training-Dynamics Diagnosis)}
Suppose a run exhibits:
\begin{itemize}[leftmargin=1.5em]
  \item training loss oscillation in first 2k steps,
  \item occasional gradient norm spikes from 12 to 140,
  \item mixed-precision overflow warnings.
\end{itemize}

Use the refresher's reasoning chain:
\begin{enumerate}[leftmargin=1.5em]
  \item \textbf{Diagnose instability source.} Early oscillation plus spikes suggests oversized effective step on high-curvature/noisy minibatches.
  \item \textbf{Apply clipping policy.} With threshold $\tau=20$, a spike at norm 140 is rescaled by factor $20/140\approx0.143$.
  \item \textbf{Check EMA startup behavior.} For $\beta_1=0.9$ and constant scalar gradient 2:
\[
m_1=0.1\cdot2=0.2,\quad m_2=0.9(0.2)+0.1(2)=0.38,\quad m_3=0.542.
\]
  Bias-corrected:
\[
\hat m_3=\frac{m_3}{1-0.9^3}=\frac{0.542}{0.271}\approx2.00.
\]
  \item \textbf{Stabilize numerics.} Increase loss-scaling headroom or use dynamic scaling rollback on overflow.
  \item \textbf{Interpret outcome.} If oscillation dampens while validation trend improves, intervention matched failure mechanism.
\end{enumerate}

\section*{Intuition Checkpoints}
\begin{itemize}[leftmargin=1.5em]
  \item Adam corrections are not magic; under drifting distributions they are approximate.
  \item Clipping controls step magnitude, not objective geometry.
  \item Mixed precision saves memory/throughput but shifts numerical failure modes.
  \item Scaling-law extrapolations degrade when data quality and objective shift.
\end{itemize}

\section*{Pointers}
\begin{itemize}[leftmargin=1.5em]
  \item Loss and KL notation: see Chapter 1 refresher.
  \item Matrix and norm geometry used by clipping and curvature intuition: see Chapter 4 refresher.
\end{itemize}

\section*{References and What To Use Them For}
\begin{itemize}[leftmargin=1.5em]
  \item \textbf{Goodfellow--Bengio--Courville, \emph{Deep Learning}} --- SGD, momentum, adaptive optimization intuition.
  \item \textbf{Kingma and Ba (Adam)} --- original adaptive-moment method and bias correction rationale.
  \item \textbf{CS229/CS231n notes} --- optimization diagnostics and practical training heuristics.
  \item \textbf{Hoffmann et al. (2022)} --- compute-optimal scaling interpretation and planning.
  \item \textbf{Bottou et al., optimization for large-scale ML} --- stochastic optimization theory background.
\end{itemize}

\section{Chapter 3 Refresher: Neural Networks}
\label{sec:math-refresher-neural-networks}

\section{Week 3 Refresher: Tokenization and Embedding Geometry}

\subsection*{Setting the Scene}
This week answers a hidden question: what mathematical object is a ``token'' once the model starts computing? The field evolved from count-based lexical models to distributed vectors, then to context-dependent embeddings in transformer systems. The critical insight was that representation is not a cosmetic preprocessing step. It defines the geometry in which similarity, retrieval, and attention operate.

\subsection*{Notation Introduced in This Chapter}
\begin{itemize}[leftmargin=1.5em]
  \item $V$: tokenizer vocabulary (finite symbol set).
  \item $\iota(\cdot)$: tokenizer mapping from text to token IDs.
  \item $\mathrm{onehot}(i)$: basis vector for token index $i$.
  \item $E$: embedding matrix.
  \item $e_i$: embedding vector for token $i$.
  \item $u,v,q,d$: generic vectors used in similarity comparisons.
\end{itemize}

\subsection*{What You Need To Recall Precisely}
\begin{itemize}[leftmargin=1.5em]
  \item \textbf{Tokenizer as function.} A tokenizer defines mapping $\iota:\text{text}\to\{1,\dots,|V|\}$. Every downstream probability is conditioned on this representation choice.
  \item \textbf{One-hot basis and embedding lookup.}
  \[
    e_i = E^\top \mathrm{onehot}(i).
  \]
  This is linear algebra row selection.
  \item \textbf{Segmentation criteria.} BPE, WordPiece, and Unigram optimize different objectives (merge frequency, likelihood-style score, probabilistic segmentation).
  \item \textbf{Similarity metrics.}
  \[
    u^\top v \quad \text{vs} \quad \cos(u,v)=\frac{u^\top v}{\|u\|\|v\|}.
  \]
  Dot product includes magnitude effects; cosine removes scale.
  \item \textbf{Anisotropy and concentration.} Contextual embeddings can cluster in narrow directional cones, weakening naive nearest-neighbor interpretations.
\end{itemize}

\subsection*{How These Results Build on Each Other}
\begin{enumerate}[leftmargin=1.5em]
  \item Tokenization defines discrete symbols.
  \item Symbol IDs become one-hot basis vectors.
  \item Embedding matrix maps one-hot vectors into continuous space.
  \item Similarity metric defines neighborhood structure.
  \item Neighborhood structure influences retrieval and downstream context selection.
\end{enumerate}

\subsection*{Reminder Glossary (Term by Term)}
\begin{itemize}[leftmargin=1.5em]
  \item \textbf{Vocabulary:} finite symbol inventory used by the model.
  \item \textbf{Token ID:} integer index assigned by tokenizer.
  \item \textbf{One-hot vector:} canonical basis vector for a token ID.
  \item \textbf{Embedding matrix:} learned lookup table from IDs to vectors.
  \item \textbf{Dot product:} magnitude-sensitive alignment score.
  \item \textbf{Cosine similarity:} direction-only alignment score.
  \item \textbf{Anisotropy:} directional imbalance in vector space.
\end{itemize}

\subsection*{Worked Example (Why Metric Choice Changes Outcomes)}
Let query vector be $q=(10,0)$ and two candidates:
\[
d_1=(1,0),\quad d_2=(1,1).
\]
Dot product ranking:
\[
q^\top d_1 = 10,\quad q^\top d_2 = 10.
\]
Tie under dot product.
Cosine ranking:
\[
\cos(q,d_1)=1,\quad \cos(q,d_2)=\frac{1}{\sqrt{2}}\approx0.707.
\]
Now $d_1$ is clearly preferred directionally.  
Reasoning implication: if your retrieval objective is semantic direction, cosine normalization is often essential; if magnitude encodes useful confidence/strength, dot-product behavior may be intentional.

\subsection*{Intuition Checkpoints}
\begin{itemize}[leftmargin=1.5em]
  \item Keep sequence-position indices separate from vocabulary indices.
  \item Tokenization choice changes sequence length and therefore perplexity comparability.
  \item Static-word-vector intuitions transfer only partially to contextual embeddings.
\end{itemize}

\subsection*{Pointers}
\begin{itemize}[leftmargin=1.5em]
  \item Probability/perplexity interpretation: see Week 1 refresher.
\end{itemize}

\subsection*{References and What To Use Them For}
\begin{itemize}[leftmargin=1.5em]
  \item \textbf{Jurafsky--Martin, \emph{SLP}} --- tokenizer design and language representation.
  \item \textbf{Strang, \emph{Linear Algebra}} --- vector spaces, norms, geometry.
  \item \textbf{Mikolov et al.} --- distributed embedding intuition.
  \item \textbf{CS224N materials} --- modern embedding/retrieval implementation context.
\end{itemize}

\section{Chapter 5 Refresher: The Transformer --- Attention as the Core Computational Primitive --- Attention, Shapes, and Complexity}

\subsection*{Setting the Scene}
Attention replaces recurrence's sequential bottleneck with explicit pairwise comparisons implemented as batched matrix multiply---linear algebra becomes architecture instead of notation garnish.

\paragraph{Step A --- bilinear compatibility.}
\textbf{Problem.} Position \(t\) must aggregate information from other positions weighted by relevance.
\textbf{Insight.} Dot products score compatibility between transformed representations (queries vs keys).
\textbf{Lineage.} Bahdanau-style attention (2015); scaled dot-product formulation crystallised in Vaswani et al.\ (2017).
\textbf{Mental model.} Think ``matching embeddings,'' not mysticism---everything reduces to tensor contractions.

\paragraph{Step B --- softmax as differentiable selector.}
\textbf{Problem.} Hard argmax blocks gradients; we need smooth convex combinations.
\textbf{Insight.} Softmax (Chapter~1 tools) turns logits into mixing weights---Preface Step~10 intuition specialised row-wise.
\textbf{Lineage.} Statistical mechanics exponential tilts $\to$ neural attention gates.
\textbf{Mental model.} Temperature inherited from decoding chapters scales sharpness identically in spirit.

\paragraph{Step C --- masks encode causal laws.}
\textbf{Problem.} Autoregressive generation forbids peeking rightward.
\textbf{Insight.} Additive \(-\infty\) constraints zero softmax mass---enforce equality constraints algebraically before nonlinear squash.
\textbf{Lineage.} Decoder masking folklore codified in Transformer stacks.
\textbf{Mental model.} Same trick handles padding tokens---different semantics, identical mathematics.

\paragraph{Step D --- quadratic surface area.}
\textbf{Problem.} Pairwise attention allocates \(T\times T\) logits per layer.
\textbf{Insight.} Long contexts explode memory/compute---motivates FlashAttention-style engineering outside this refresher's scope but inside systems reasoning.
\textbf{Lineage.} Dao et al.\ IO-aware kernels (2022 era).
\textbf{Mental model.} Preface Step~6 variance arguments resurface when chunking randomises attention pathways.

If shapes align, attention is honest matrix cookbookery---when shapes drift, bugs hide silently.

\subsection*{Notation Introduced in This Chapter}
\begin{itemize}[leftmargin=1.5em]
  \item $T$: sequence length.
  \item $d_{\text{model}}$: model hidden width.
  \item $H$: number of attention heads.
  \item $d_k$: per-head key/query dimension.
  \item $Q,K,V$: query/key/value matrices.
  \item $M$: additive mask matrix applied before softmax.
\end{itemize}

\subsection*{What You Need To Recall Precisely}
\begin{itemize}[leftmargin=1.5em]
  \item \textbf{Scaled attention:}
  \[
    \mathrm{Attn}(Q,K,V)=\mathrm{softmax}\!\left(\frac{QK^\top+M}{\sqrt{d_k}}\right)V.
  \]
  Read row-wise: row $t$ of $QK^\top$ lists dot-product similarities between position-$t$ query vectors and every key vector; softmax turns those logits into nonnegative weights summing to $1$; multiplying those weights into $V$ mixes value vectors---attention outputs are convex combinations of values attended-to.

  \textbf{Why divide by $\sqrt{d_k}$?}
  If components of queries and keys are $\mathcal{O}(1)$ with roughly independent variation across $d_k$ dimensions, dot-product magnitudes grow \(\Theta(\sqrt{d_k})\) in expectation (length scales like Euclidean norm). Without scaling, softmax logits become sharply peaked as width grows, driving gradients toward sparse argmax-like behaviour and harming trainability. Rescaling keeps typical logits \(\mathcal{O}(1)\) across widths.

  \item \textbf{Mask semantics:} causal mask forbids looking right; padding mask blocks non-content locations.
  \item \textbf{Head decomposition:} independent projections create multiple interaction subspaces.
  \item \textbf{Residual and normalization:} residual streams pass representations forward as \(x\leftarrow x+\mathrm{Block}(x)\), letting blocks learn corrections around identity rather than full maps from scratch. Normalisation (LayerNorm/RMSNorm inside Transformer blocks) recentres/rescales activations per token or position so dot products do not drift orders of magnitude mid-training---pairs with residual paths to control gradient scale through depth.
  \item \textbf{Complexity source:} score matrix is $T\times T$, so interaction cost grows with sequence length.
\end{itemize}

\subsection*{How These Results Build on Each Other}
\begin{enumerate}[leftmargin=1.5em]
  \item Project input vectors into $Q,K,V$.
  \item Form pairwise compatibility scores via $QK^\top$.
  \item Apply structural constraints with mask $M$.
  \item Normalize row-wise with softmax.
  \item Aggregate values and return contextualized representations.
\end{enumerate}

\subsection*{Reminder Glossary (Term by Term)}
\begin{itemize}[leftmargin=1.5em]
  \item \textbf{Query:} representation of what current position is asking for.
  \item \textbf{Key:} representation of what each position offers.
  \item \textbf{Value:} content payload mixed by attention weights.
  \item \textbf{Mask:} hard structural restriction on allowable dependencies.
  \item \textbf{Head:} one parallel attention subspace.
  \item \textbf{KV cache:} stored keys/values for efficient autoregressive decoding.
\end{itemize}

\subsection*{Worked Mini Example (Shape Sanity)}
Let $T=4$, $d_{\text{model}}=8$, heads $H=2$, so $d_k=4$.
\[
Q,K,V \in \mathbb{R}^{4\times 4}\ \text{per head},\quad QK^\top \in \mathbb{R}^{4\times 4}.
\]
After softmax over rows, multiply by $V$ to return $\mathbb{R}^{4\times 4}$ per head, concatenate to $\mathbb{R}^{4\times 8}$.

\subsection*{Worked Example (Reasoning Chain with Causal Mask)}
For token position $t=3$ in a length-4 causal sequence, allowed keys are positions $\{1,2,3\}$ only.  
In row 3 of $QK^\top$, mask sets column 4 score to $-\infty$ before softmax.  
Result: attention weight at $(3,4)$ is exactly 0.  
Interpretation: causality is enforced at score level, not corrected later at output level.

\subsection*{Intuition Checkpoints}
\begin{itemize}[leftmargin=1.5em]
  \item ``Masking'' is a constraint on attention scores, not a post-hoc output edit.
  \item Quadratic attention cost refers to score interactions, not all block components equally.
  \item KV cache changes decoding-time behavior, not pretraining-time complexity class.
\end{itemize}

\subsection*{Further Refresh}
\begin{itemize}[leftmargin=1.5em]
  \item Shape and vector geometry reminders: see Chapter 4 refresher.
  \item Softmax stability reminders: see Chapter 1 refresher.
  \item \textbf{Vaswani et al. (2017)} --- architectural equations and design rationale.
  \item \textbf{Annotated Transformer} --- equation-by-equation implementation mapping.
\end{itemize}

\chapter{Chapter 6 Refresher: Pretraining, GPT-Style Models, and In-Context Learning --- Pretraining Objective and Decoding Controls}

\section*{Setting the Scene}
Training optimises Preface Steps~8--10 on teacher-forced contexts; deployment samples from the learned conditional family under budgets and safety envelopes---same \(\theta\), different \emph{pushforward} map from logits to tokens.

\paragraph{Step A --- Autoregressive likelihood reminder.}
\textbf{Problem.} Guarantee coherence across positions while fitting data.
\textbf{Insight.} Factorisation \(\prod_t p_\theta(x_t\mid x_{<t})\) ties optimisation to causal graphs already rehearsed in Preface Step~3.
\textbf{Lineage.} Shannon's entropy estimates for English $\to$ modern causal transformers.
\textbf{Mental model.} Teacher forcing mismatches free-running drift---monitor both regimes.

\paragraph{Step B --- Inference-time reshaping.}
\textbf{Problem.} Raw softmax may be too peaked or too flat for products.
\textbf{Insight.} Temperature, top-\(k\), top-\(p\) apply measure transforms \emph{after} forward pass---no gradient necessarily flows (pure inference policy).
\textbf{Lineage.} Nucleus sampling (Holtzman et al.) formalises mass truncation psychology.
\textbf{Mental model.} Treat decoding knobs as UI levers on Shannon entropy (Preface Step~9).

\paragraph{Step C --- Scaling budgets.}
\textbf{Problem.} Engineers trade parameters vs tokens under compute roofs.
\textbf{Insight.} Empirical iso-loss laws operationalise Preface Step~7 qualitatively at trillion-token scales.
\textbf{Lineage.} Kaplan (2020); Chinchilla (2022).
\textbf{Mental model.} Laws fit historical pipelines---extrapolate cautiously when data composition shifts.

Decoding is theatre lighting on the same actor (\(\theta\)); training chooses the script likelihood, decoding chooses spotlight intensity.

\section*{Notation Introduced in This Chapter}
\begin{itemize}[leftmargin=1.5em]
  \item $z_i$: logit for candidate token $i$.
  \item $p_i$: probability of candidate token $i$ after normalization.
  \item $\tau$: temperature parameter.
  \item $k$: top-$k$ truncation count.
  \item $p$: top-$p$ cumulative mass threshold.
  \item $N,D$: model/data scale variables in scaling discussions.
\end{itemize}

\section*{What You Need To Recall Precisely}
\begin{itemize}[leftmargin=1.5em]
  \item \textbf{Autoregressive NLL objective.}
  For corpus sequences \(x_{1:T}\), training maximises \(\sum_t \log p_\theta(x_t\mid x_{<t})\), equivalently minimises per-token cross-entropy averaged over positions and documents. This is the Chapter~1 chain rule factorisation evaluated inside the transformer stack that produces conditional softmax parameters at each \(t\).
  \item \textbf{Temperature transform:}
  \[
    p_i(\tau)=\frac{\exp(z_i/\tau)}{\sum_j \exp(z_j/\tau)}.
  \]
  \item \textbf{Top-$k$ and top-$p$ (ordered operations).}
  Start from full softmax masses \(p_i\) over vocabulary.
  \begin{enumerate}[leftmargin=1.5em]
    \item \textbf{Top-$k$:} keep the $k$ largest probabilities, zero the rest, divide remaining masses by their sum so they again form a distribution on the surviving support.
    \item \textbf{Top-$p$ (nucleus):} sort probabilities descending, greedily add tokens until cumulative mass reaches threshold $p\in(0,1]$, zero outside that prefix, renormalise on the kept set.
  \end{enumerate}
  Both policies intentionally discard probability mass outside a subset of vocabulary entries before sampling; they reshape diversity without changing trained weights.
  \item \textbf{Objective/metric distinction:} training objective may not align perfectly with deployment metric.
  \item \textbf{Scaling fits:} compute/data/parameter fronts are empirical and objective-dependent.
\end{itemize}

\section*{Pointers}
\begin{itemize}[leftmargin=1.5em]
  \item CE/KL/perplexity: see Chapter 1 refresher.
  \item Optimizer and scaling-law mechanics: see Chapter 2 refresher.
  \item Attention architecture assumptions: see Chapter 5 refresher.
\end{itemize}

\section*{How These Results Build on Each Other}
\begin{enumerate}[leftmargin=1.5em]
  \item Train model to minimize autoregressive NLL.
  \item Convert logits to probabilities at inference.
  \item Apply decoding transform (temperature/truncation).
  \item Sample outputs from transformed distribution.
  \item Evaluate behavior with metrics that may differ from training objective.
\end{enumerate}

\section*{Reminder Glossary (Term by Term)}
\begin{itemize}[leftmargin=1.5em]
  \item \textbf{Autoregressive objective:} predict each token conditioned on prefix.
  \item \textbf{Decoding policy:} rule for turning model distribution into emitted tokens.
  \item \textbf{Temperature:} entropy-control knob on logits.
  \item \textbf{Top-$k$:} fixed-cardinality truncation.
  \item \textbf{Top-$p$:} mass-based truncation.
  \item \textbf{Scaling law:} empirical relation between loss and resource axes.
\end{itemize}

\section*{Worked Mini Example (Temperature)}
Take logits $z=(2,1,0)$.
With $\tau=1$:
\[
p\approx(0.665,0.245,0.090).
\]
With $\tau=2$:
\[
p\approx(0.506,0.307,0.186).
\]
Higher temperature flattens the distribution, increasing diversity and risk simultaneously.

\section*{Worked Example (Decode Policy Choice for Same Model)}
Suppose next-token distribution after softmax is:
\[
(0.52,0.24,0.12,0.07,0.05).
\]
Under top-$k$ with $k=2$, support becomes first two tokens, renormalized to
\[
\left(\frac{0.52}{0.76},\frac{0.24}{0.76}\right)\approx(0.684,0.316).
\]
Under top-$p$ with $p=0.9$, first four tokens are retained (mass $0.95$), renormalized to:
\[
(0.547,0.253,0.126,0.074).
\]
Same trained model, different generation behavior due only to policy.

\section*{Intuition Checkpoints}
\begin{itemize}[leftmargin=1.5em]
  \item Decoding does not retrain the model; it samples from a transformed output distribution.
  \item Top-$k$ and top-$p$ can behave similarly in peaked distributions and very differently in flat distributions.
  \item Apparent ``emergence'' can be caused by metric discontinuities, not abrupt capability phase transitions.
\end{itemize}

\section*{Further Refresh}
\begin{itemize}[leftmargin=1.5em]
  \item \textbf{Hoffmann et al. (2022)} --- compute-optimal scaling interpretation.
  \item \textbf{Holtzman et al.} --- nucleus sampling and degeneration analysis.
  \item Standard decoding notes from CS224N/modern LLM systems courses.
\end{itemize}

\chapter{Chapter 7 Refresher: Efficiency, Scaling, and Mixture-of-Experts --- Systems Arithmetic and Sparse Routing}

\section*{Setting the Scene}
Once probabilities and networks are defined, training becomes physics on silicon: bytes move, joules dissipate, routers stall. Preface Step~6 (variance of sums) quietly lurks behind contention noise.

\paragraph{Step A --- Roofline worldview.}
\textbf{Problem.} MFLOPS dashboards lie when kernels stream weights constantly.
\textbf{Insight.} Compare arithmetic intensity (FLOPs/byte) against bandwidth-induced slopes---whichever bound bites first wins.
\textbf{Lineage.} Williams/Waters/Patterson roofline model (2009) repurposed across ML hardware discourse.
\textbf{Mental model.} If intensity $<$ ridge crossover, buying more TFLOPs wastes money---buy bandwidth or locality instead.

\paragraph{Step B --- Sparsity swaps villains.}
\textbf{Problem.} Dense matrices saturate compute; sparse routing saves FLOPs but reroutes traffic.
\textbf{Insight.} MoE trades matmuls for all-to-all collectives---new bottleneck on wires (Preface Step~6 covariance spikes across devices).
\textbf{Lineage.} Shazeer's Switch Transformer et seq.
\textbf{Mental model.} Plot effective \(\lambda/\mu\) including dispatch latency before celebrating parameter counts.

\paragraph{Step C --- Capacity factor as queue abstraction.}
\textbf{Problem.} Experts overload like servers under burst arrivals.
\textbf{Insight.} Capacity factor overprovisions slots analogous to multi-server queue stability margins---drops mirror admission control.
\textbf{Lineage.} Queueing tradition (Kendall notation) meets routing papers.
\textbf{Mental model.} Tail latency dominates UX---average FLOPs/token misleads.

Systems graphs reinterpret probability routines as throughput limits; migrating bottlenecks obey no moral ranking beyond measurement.

\section*{Notation Introduced in This Chapter}
\begin{itemize}[leftmargin=1.5em]
  \item $I$: arithmetic intensity (FLOPs per byte).
  \item $k$: number of selected experts per token in top-$k$ routing.
  \item $N_{\text{active}}$: active parameters per token.
  \item $N_{\text{total}}$: total resident parameters.
  \item CF: capacity factor (expert slot overprovision multiplier).
\end{itemize}

\section*{What You Need To Recall Precisely}
\begin{itemize}[leftmargin=1.5em]
  \item \textbf{Bottleneck triad:} compute, memory bandwidth, and interconnect communication.
  \item \textbf{MoE routing:} compute gate logits per token (router network); softmax yields mixture weights; keep top-$k$ experts, zero or mask others, renormalise weights on the surviving experts so each token still mixes convexly over experts it dispatches to; aggregate expert outputs with those weights. Routing induces sparse activation patterns separate from parameter residency.
  \item \textbf{Capacity factor:} finite per-expert token slots induce drop/reroute decisions.
  \item \textbf{Arithmetic intensity:} $\text{FLOPs}/\text{byte}$ predicts memory-bound vs compute-bound behavior.
  \item \textbf{Active vs total parameters:} sparse models can have large total memory footprint while using small active compute per token.
\end{itemize}

\section*{How These Results Build on Each Other}
\begin{enumerate}[leftmargin=1.5em]
  \item Compute workload demand (FLOPs and bytes/token).
  \item Compare with device and interconnect limits.
  \item Identify bottleneck (compute, memory, communication).
  \item Apply sparse routing and reassess new bottleneck.
  \item Validate with throughput/latency measurements.
\end{enumerate}

\section*{Reminder Glossary (Term by Term)}
\begin{itemize}[leftmargin=1.5em]
  \item \textbf{Arithmetic intensity:} FLOPs per byte transferred.
  \item \textbf{Dispatch/combine:} route tokens to experts and gather outputs.
  \item \textbf{Capacity factor:} overprovisioning multiplier for expert token slots.
  \item \textbf{Token drop:} overflow handling when expert capacity exceeded.
  \item \textbf{All-to-all:} communication primitive central to MoE scaling.
\end{itemize}

\section*{Worked Mini Example (Intensity Intuition)}
Assume a kernel performs $2\times10^{12}$ FLOPs and moves $8\times10^{11}$ bytes.
\[
I=\frac{2\times10^{12}}{8\times10^{11}}=2.5\ \text{FLOPs/byte}.
\]
If hardware requires much higher intensity for compute saturation, this kernel is memory-bound.

\section*{Worked Example (MoE Bottleneck Migration)}
Assume dense model requires 100 units compute and 40 units communication per batch.  
After MoE sparsification, active compute drops to 45, but dispatch/combine communication rises to 80.  
System shifts from compute-bound to communication-bound.  
Interpretation: sparse compute improved one axis but exposed another, so optimization must target network scheduling or expert placement.

\section*{Intuition Checkpoints}
\begin{itemize}[leftmargin=1.5em]
  \item Sparse compute does not imply sparse communication.
  \item MoE can reduce compute/token while increasing routing complexity.
  \item High model parameter count can coexist with low active compute and high memory residency.
\end{itemize}

\section*{Further Refresh}
\begin{itemize}[leftmargin=1.5em]
  \item Optimization and memory accounting: see Chapter 2 refresher.
  \item Attention complexity and KV behavior: see Chapter 5 refresher.
  \item Decoding/token-rate implications: see Chapter 6 refresher.
  \item \textbf{Shazeer et al. / Switch Transformer papers} --- sparse routing mechanics and load balancing.
  \item \textbf{Harchol-Balter performance modeling} --- bottleneck and throughput reasoning.
\end{itemize}

\chapter{Chapter 8 Refresher: Model Adaptation, LoRA, Quantization, and Deployment Reality --- Low-Rank Adaptation and Quantization}

\section*{Setting the Scene}
Deployment freezes most of \(\theta\) learned via Preface Step~8-style objectives and tunes cheap corrections under storage/power envelopes---approximation theory meets accounting.

\paragraph{Step A --- Low-rank hypothesis.}
\textbf{Problem.} Full matrices encode redundancy; updating all entries per task is wasteful.
\textbf{Insight.} SVD decomposes any matrix into a sum of rank-1 outer products ordered by singular value magnitude. Restrict \(\Delta W\) to rank-\(r\) manifolds---implicit claim task signals concentrate in few directions (Preface Step~6 thinks about singular values informally).
\textbf{Lineage.} Eckart--Young (1936) best low-rank approximation theorem; Hu et al.\ LoRA (2021) popularised modern adapter lore.
\textbf{Mental model.} Monitor singular spectrum drift when skeptics question rank budgets---a fast-decaying spectrum justifies small $r$; a flat spectrum does not.

\paragraph{Step B --- Quantisation as noisy channel.}
\textbf{Problem.} Continuous weights exceed SRAM budgets.
\textbf{Insight.} Rounding injects controlled distortion reminiscent of quantised communication channels---error scales with dynamic range misuse.
\textbf{Lineage.} Dettmers QLoRA et al.\ marry low-rank updates with int4 storage.
\textbf{Mental model.} Outliers behave like heavy-tail Preface Step~6 violations---scale grids adapt or suffer.

\paragraph{Step C --- Gate on measured drift.}
\textbf{Problem.} Approximations trade compute for risk.
\textbf{Insight.} Acceptance tests slice datasets like calibration bins---aggregate scores hide subgroup collapses tied to Preface Step~5 expectations.

Treat every shortcut as falsifiable: publish metrics or admit folklore.

\section*{Notation Introduced in This Chapter}
\begin{itemize}[leftmargin=1.5em]
  \item $W$: base weight matrix, $\Delta W$: adaptation update.
  \item $A,B$: low-rank factors in LoRA update.
  \item $d,k$: layer input/output dimensions.
  \item $r$: LoRA rank.
  \item $\alpha$: LoRA scaling constant.
  \item $\rho$: trainable parameter fraction versus full fine-tuning.
  \item $U,\Sigma,V$: left singular vectors, diagonal singular values, right singular vectors in SVD.
  \item $\sigma_i$: the $i$-th singular value (conventionally $\sigma_1 \ge \sigma_2 \ge \cdots \ge 0$).
  \item $\|M\|_F$: Frobenius norm, $\sqrt{\sum_{i,j}M_{ij}^2} = \sqrt{\sum_i \sigma_i^2}$.
\end{itemize}

\section*{What You Need To Recall Precisely}
\begin{itemize}[leftmargin=1.5em]
  \item \textbf{Singular Value Decomposition (SVD) and the low-rank approximation theorem.}

  Every real matrix $M \in \mathbb{R}^{d \times k}$ factors as
  \[
    M = U\Sigma V^\top,
  \]
  where $U \in \mathbb{R}^{d\times d}$ and $V \in \mathbb{R}^{k\times k}$ are orthogonal, and $\Sigma$ is diagonal with non-negative entries $\sigma_1 \ge \sigma_2 \ge \cdots \ge 0$. Expanding the product gives a sum of rank-1 outer products:
  \[
    M = \sum_{i=1}^{\min(d,k)} \sigma_i\, \mathbf{u}_i \mathbf{v}_i^\top.
  \]
  \emph{Eckart--Young theorem.} For any target rank $r$, the best rank-$r$ approximation in Frobenius norm (and in spectral norm) is obtained by keeping the top $r$ terms:
  \[
    M_r = \sum_{i=1}^{r} \sigma_i\, \mathbf{u}_i \mathbf{v}_i^\top,
    \qquad
    \|M - M_r\|_F = \sqrt{\sigma_{r+1}^2 + \cdots + \sigma_{\min(d,k)}^2}.
  \]
  No other rank-$r$ matrix is closer to $M$ in Frobenius norm. This is the theoretical foundation for using low-rank matrices as approximations.

  \emph{Connection to LoRA.} The LoRA hypothesis is that the task-relevant update $\Delta W$ has a rapidly decaying singular spectrum---most of its ``signal'' lives in a low-dimensional subspace. Restricting $\Delta W = \frac{\alpha}{r}BA$ (with $B \in \mathbb{R}^{d\times r}$, $A \in \mathbb{R}^{r\times k}$) is equivalent to parameterising the rank-$r$ subspace directly. Eckart--Young justifies this when the true update's singular values drop off quickly; when the spectrum is flat, a small $r$ will miss significant directions and underfitting follows.

  \emph{Diagnosing rank budget.} Inspect the singular value spectrum of $\Delta W$ after a full fine-tuning run (if affordable as a baseline). Rapid decay --- first few $\sigma_i$ capturing $\ge 90\%$ of $\|{\Delta W}\|_F$ --- supports a small LoRA rank. Slow decay flags either a hard task shift or a dataset with diverse sub-distributions, both arguing for larger $r$ or continued pretraining.

  \item \textbf{LoRA update form and intuition.}
  \[
    \Delta W = \frac{\alpha}{r}BA,\quad B\in\mathbb{R}^{d\times r},\ A\in\mathbb{R}^{r\times k}.
  \]
  Rank \(r\ll \min(d,k)\) forces \(\Delta W\) to live in an \(r\)-dimensional subspace of the full \(d\times k\) parameter space; optimisation restricts plasticity to directions where adaptation concentrates empirically, exchanging expressiveness for parameter savings.

  \item \textbf{Trainable parameter fraction:}
  \[
    \rho=\frac{r(d+k)}{dk}.
  \]
  \item \textbf{Quantization model:} map tensors through affine rescaling (per-channel or per-tensor scale/zero-point) into finite-bit grids---often asymmetric ranges capture outliers better than naive symmetric buckets. Quantisation noise concentrates where scales mismatch empirical tails (especially outliers), triggering subgroup drift unrelated to aggregate perplexity.
  \item \textbf{Error tradeoff:} memory/bandwidth gains versus approximation and calibration drift risk.
  \item \textbf{Budget accounting:} adaptation and base-model memory are separate components.
\end{itemize}

\section*{How These Results Build on Each Other}
\begin{enumerate}[leftmargin=1.5em]
  \item SVD decomposes any matrix into rank-1 components ordered by importance.
  \item Eckart--Young establishes that truncating to rank $r$ is the optimal approximation under Frobenius/spectral norm.
  \item The low-rank hypothesis applies this to weight updates: restrict $\Delta W$ to rank $r$ if task signal concentrates in few directions.
  \item LoRA parameterises the rank-$r$ subspace explicitly with factors $B, A$.
  \item Quantization adds a second approximation layer (bit-width reduction) with independent noise characteristics.
  \item Deployment gating measures drift from both approximations before release.
\end{enumerate}

\section*{Reminder Glossary (Term by Term)}
\begin{itemize}[leftmargin=1.5em]
  \item \textbf{Singular value decomposition (SVD):} factorisation $M = U\Sigma V^\top$ into orthogonal factors and a diagonal matrix of singular values; every real matrix has one.
  \item \textbf{Singular value ($\sigma_i$):} non-negative scalar measuring the ``size'' of the $i$-th rank-1 component; $\sigma_1 \ge \sigma_2 \ge \cdots$.
  \item \textbf{Singular spectrum:} the sequence $(\sigma_1, \sigma_2, \ldots)$; fast decay supports low-rank approximation, flat spectrum does not.
  \item \textbf{Eckart--Young theorem:} the best rank-$r$ approximation in Frobenius (or spectral) norm is the truncated SVD $M_r$; there is no better rank-$r$ matrix.
  \item \textbf{Frobenius norm:} $\|M\|_F = \sqrt{\sum_{i,j}M_{ij}^2} = \sqrt{\sum_i \sigma_i^2}$; natural measure of total approximation error.
  \item \textbf{Rank:} dimensionality of adaptation subspace; in LoRA, the number of singular directions parameterised explicitly.
  \item \textbf{LoRA scaling factor:} controls update magnitude relative to rank.
  \item \textbf{Quantizer:} mapping from continuous values to discrete levels.
  \item \textbf{Calibration set:} data used to estimate quantization scales.
  \item \textbf{Drift threshold:} maximum acceptable quality loss for deployment.
\end{itemize}

\section*{Worked Mini Example (SVD and Eckart--Young)}
Let $M = \begin{pmatrix}3&0\\0&1\\0&0\end{pmatrix}$. The SVD is $M = U\Sigma V^\top$ with
\[
  U = I_3,\quad \Sigma = \begin{pmatrix}3&0\\0&1\\0&0\end{pmatrix},\quad V = I_2,
\]
giving singular values $\sigma_1=3$, $\sigma_2=1$. The best rank-1 approximation keeps only $\sigma_1$:
\[
  M_1 = \begin{pmatrix}3\\0\\0\end{pmatrix}(1\ 0) = \begin{pmatrix}3&0\\0&0\\0&0\end{pmatrix},
  \qquad
  \|M - M_1\|_F = \sigma_2 = 1.
\]
No other rank-1 matrix achieves Frobenius error below 1. The ratio $\sigma_1^2/(\sigma_1^2+\sigma_2^2) = 9/10 = 90\%$ of total squared norm is captured by rank 1---a fast-decaying spectrum that supports the low-rank approximation.

\section*{Worked Mini Example (Parameter Fraction)}
Let $d=k=4096$ and $r=16$:
\[
\rho=\frac{16(4096+4096)}{4096\cdot4096}
=\frac{131072}{16777216}\approx 0.0078.
\]
This is about $0.78\%$ of full fine-tuning parameters for that layer.

\section*{Worked Example (Quality Gate Under Approximation)}
Suppose full fine-tuned baseline scores 78.4 on target benchmark.  
LoRA+quantized candidate scores 77.1 overall, but subgroup analysis shows:
\begin{itemize}[leftmargin=1.5em]
  \item common cases: -0.6
  \item rare long-context cases: -3.8
\end{itemize}
Reasoning: average drop may appear acceptable, but concentrated degradation violates deployment intent for long-context tasks.  
Decision: reject current quantization profile or introduce subgroup-aware gate.

\section*{Intuition Checkpoints}
\begin{itemize}[leftmargin=1.5em]
  \item SVD always exists and is unique (up to sign); it is not an approximation --- it is an exact factorisation. The approximation step is the truncation to rank $r$.
  \item Eckart--Young gives a lower bound on approximation error: $\|M-M_r\|_F \ge \sigma_{r+1}$. No representation trick can do better for a given rank.
  \item A low-rank update is a hypothesis about task-adaptation dimensionality, not a theorem. Inspect the singular spectrum to validate it.
  \item Quantization quality is distribution-dependent; outlier channels can dominate failure.
  \item ``Near parity'' must be defined metric-by-metric, not as a vague label.
\end{itemize}

\section*{Further Refresh}
\begin{itemize}[leftmargin=1.5em]
  \item Rank and linear-algebra background: see Chapter 4 refresher.
  \item Optimizer/memory arithmetic: see Chapter 2 refresher.
  \item Systems bottlenecks and throughput: see Chapter 7 refresher.
  \item \textbf{LoRA (Hu et al.)} --- low-rank adaptation mechanics.
  \item \textbf{QLoRA (Dettmers et al.)} --- quantization plus adapter training.
\end{itemize}

\chapter{Chapter 9 Refresher: Alignment, Instruction Tuning, and Preference Optimization --- Preferences, KL Regularization, and Policy Updates}

\scenelabel
Alignment reframes Preface Steps~2--4: instead of conditioning only on observed tokens, we condition on human comparisons and regularise policies toward trusted references.

\paragraph{Step A --- Pairwise likelihood.}
\textbf{Problem.} Scalar correctness insufficient; humans supply partial orders.
\textbf{Insight.} Bradley--Terry models lift score gaps into Bernoulli probabilities---Preface Step~2 specialised to preference events.
\textbf{Lineage.} Thurstone/B Bradley--Terry tradition; Rafailov et al.\ DPO (2023) streamlines derivatives.
\textbf{Mental model.} Relative labels estimate latent utilities; absolute calibration remains optional extra.

\paragraph{Step B --- KL anchor.}
\textbf{Problem.} Unconstrained policy optimisation hallucinates unsupported modes.
\textbf{Insight.} KL penalties replay Preface Step~10 mismatch narrative between \(\pi_\theta\) and \(\pi_{\mathrm{ref}}\).
\textbf{Lineage.} Schulman/PPO lineage intersects RLHF practice (Stiennon, Ouyang, subsequent stacks).
\textbf{Mental model.} \(\beta\) trades helpfulness vs behavioural drift---no universal optimum without governance context.

\paragraph{Step C --- Support hazards.}
\textbf{Problem.} Zero reference probability yields exploding log-ratios.
\textbf{Insight.} Conditioning events need positive measure---same vigilance as Preface Step~3 warned for chain factors.

Preference optimisation is probabilistic inference under sociotechnical data---mathematics unchanged, likelihood rewritten.

\section*{Notation Introduced in This Chapter}
\begin{itemize}[leftmargin=1.5em]
  \item $x$: prompt/context.
  \item $y^+, y^-$: preferred and dispreferred responses.
  \item $r(x,y)$: reward/score function.
  \item $\sigma(\cdot)$: logistic sigmoid function.
  \item $\pi$: current policy distribution.
  \item $\pi_{\mathrm{ref}}$: reference policy.
  \item $\beta$: KL regularization coefficient.
\end{itemize}

\section*{What You Need To Recall Precisely}
\begin{itemize}[leftmargin=1.5em]
  \item \textbf{Bradley--Terry form:}
  \[
    P(y^+ \succ y^- \mid x)=\sigma(r(x,y^+)-r(x,y^-)).
  \]
  Write margin \(\Delta r = r(x,y^+)-r(x,y^-)\). In Bradley--Terry style models this acts like a log-odds difference: \(\sigma(\Delta r)\) matches how humans choose stochastically under utility noise; pairwise losses therefore encourage reward gaps consistent with observed preferences.

  \item \textbf{KL regularization role.}
  Optimising preferences alone can peel \(\pi_\theta\) away from pretrained \(\pi_{\mathrm{ref}}\) into unsupported linguistic regimes where rewards mis-generalise. Adding \(\beta\, D_{\mathrm{KL}}(\pi_\theta \,\|\, \pi_{\mathrm{ref}})\) (or sampled surrogates thereof) penalises drift measured against data-supported behaviour while pairwise margins tilt \(\pi_\theta\) toward helpful completions.

  \item \textbf{Pairwise objective geometry:} score differences (margins) drive learning signal.
  \item \textbf{Support condition:} near-zero reference probability destabilizes log-ratio terms.
  \item \textbf{Objective-vs-behavior distinction:} optimizing preference loss does not automatically guarantee safety or honesty on all distributions.
\end{itemize}

\section*{How These Results Build on Each Other}
\begin{enumerate}[leftmargin=1.5em]
  \item Convert pairwise labels into likelihood terms.
  \item Learn scoring or policy update from pairwise margins.
  \item Add KL regularization to bound policy drift.
  \item Validate on held-out safety/helpfulness metrics.
\end{enumerate}

\section*{Reminder Glossary (Term by Term)}
\begin{itemize}[leftmargin=1.5em]
  \item \textbf{Preference pair:} chosen vs rejected response for same prompt.
  \item \textbf{Margin:} score difference between preferred and non-preferred outputs.
  \item \textbf{KL penalty:} regularizer limiting divergence from reference policy.
  \item \textbf{Support mismatch:} missing probability mass causing unstable ratios.
  \item \textbf{Over-optimization:} objective improves while true quality degrades.
\end{itemize}

\section*{Worked Mini Example (Pairwise Probability)}
If $r(x,y^+)=2.0$ and $r(x,y^-)=0.5$:
\[
P(y^+\succ y^-|x)=\sigma(1.5)\approx0.817.
\]
Interpretation: the model assigns roughly $82\%$ preference probability to $y^+$ under this reward gap.

\section*{Worked Example (Preference Gain vs Safety Drift)}
Suppose after tuning:
\begin{itemize}[leftmargin=1.5em]
  \item preference win-rate improves from 62\% to 73\%,
  \item refusal-on-unsafe-prompts worsens from 91\% to 84\%.
\end{itemize}
Reasoning: preference objective improved but safety behavior regressed.  
Implication: keep preference metric and safety metric as separate acceptance gates; one does not substitute for the other.

\section*{Intuition Checkpoints}
\begin{itemize}[leftmargin=1.5em]
  \item Lower KL does not mean better helpfulness; it means less policy drift.
  \item Better reward-model fit does not guarantee lower downstream harm.
  \item Pairwise preferences encode relative order, not absolute truth.
\end{itemize}

\section*{Further Refresh}
\begin{itemize}[leftmargin=1.5em]
  \item KL/cross-entropy foundations: see Chapter 1 refresher.
  \item Optimization mechanics: see Chapter 2 refresher.
  \item Adaptation substrate (LoRA/QLoRA): see Chapter 8 refresher.
  \item \textbf{Christiano/Stiennon/Ouyang papers} --- RLHF workflow and assumptions.
  \item \textbf{Rafailov et al. (DPO)} --- direct preference objective derivation.
\end{itemize}

\chapter{Chapter 10 Refresher: LLM Systems --- Retrieval, Tool Use, and Evaluation Harnesses --- Retrieval Metrics and Evaluation Arithmetic}

\section*{Setting the Scene}
Retrieval-augmented answering concatenates two stochastic modules: recall relevant evidence (set-valued event), then condition generation on that evidence (Preface Step~2 again). One scalar cannot isolate which module failed.

\paragraph{Step A --- Set-valued relevance.}
\textbf{Problem.} Truth may admit multiple supporting documents.
\textbf{Insight.} Precision/recall generalise Preface Step~1 counting logic to partial overlaps between retrieved and gold sets.
\textbf{Lineage.} Classical IR (Manning et al.); neural hybrids layer embedding geometry from Chapter~4.

\paragraph{Step B --- Conditional generation quality.}
\textbf{Problem.} Good retriever + weak synthesiser still ships nonsense.
\textbf{Insight.} Evaluate \(\mathcal{L}\) or task metrics \emph{given} oracle contexts vs live contexts---difference isolates grounding gaps.

\paragraph{Step C --- Faithfulness as constrained decoding.}
\textbf{Problem.} Fluent generators invent citations.
\textbf{Insight.} Faithfulness metrics approximate hard logical entailment with softer overlap checks---engineering surrogate for impossible full semantics.

Decomposition honors Preface Step~5 linearity instinct: diagnose components before averaging them blindly.

\section*{Notation Introduced in This Chapter}
\begin{itemize}[leftmargin=1.5em]
  \item $k$: retrieval cutoff.
  \item $R_k$: top-$k$ retrieved set.
  \item $G$: gold relevant set.
  \item Precision@$k$, Recall@$k$: standard ranking metrics.
  \item $q$: query embedding/vector.
  \item $d$: document/chunk embedding/vector.
\end{itemize}

\section*{What You Need To Recall Precisely}
\begin{itemize}[leftmargin=1.5em]
  \item \textbf{Ranking metrics.}
  Let \(R_k\) be top-$k$ retrieved IDs and \(G\) gold relevant IDs with $|G|$ possibly larger than $k$.
  \[
    \mathrm{Prec@}k=\frac{|R_k \cap G|}{k},\qquad
    \mathrm{Rec@}k=\frac{|R_k \cap G|}{|G|}.
  \]
  Precision stresses cleanliness of short rankings; recall stresses coverage of all positives---trade-offs arise whenever $|G|$ is large relative to $k$. Mean reciprocal rank (MRR) averages \(1/\text{rank of first relevant}\) across queries---rewarding systems that surface \emph{one} correct doc quickly rather than maximising exhaustive recall.
  \item \textbf{Sparse vs dense retrieval:} lexical overlap and embedding similarity capture different relevance modes.
  \item \textbf{Fusion methods:} rank fusion stabilizes retrieval when signals disagree.
  \item \textbf{Evaluation decomposition:} retrieval-stage errors and generation-stage errors require separate diagnostics.
  \item \textbf{Grounding metrics:} answer support must be checked against retrieved evidence, not linguistic plausibility.
\end{itemize}

\section*{How These Results Build on Each Other}
\begin{enumerate}[leftmargin=1.5em]
  \item Define relevance and evidence assumptions.
  \item Measure retrieval coverage/precision.
  \item Measure generation quality conditioned on retrieved context.
  \item Measure grounding/faithfulness of final claims.
  \item Localize failure stage before proposing fix.
\end{enumerate}

\section*{Reminder Glossary (Term by Term)}
\begin{itemize}[leftmargin=1.5em]
  \item \textbf{Gold relevance set:} trusted set of truly relevant items.
  \item \textbf{Precision@k:} fraction of top-$k$ retrieved items that are relevant.
  \item \textbf{Recall@k:} fraction of relevant items captured in top-$k$.
  \item \textbf{Reranker:} second-stage scoring model for refined ranking.
  \item \textbf{Faithfulness:} proportion of answer claims supported by context.
\end{itemize}

\section*{Worked Mini Example (Precision and Recall@k)}
Suppose the gold relevant set has 4 items. Top-5 retrieval returns 3 relevant items.
\[
\mathrm{Precision@5}=\frac{3}{5}=0.6,\qquad
\mathrm{Recall@5}=\frac{3}{4}=0.75.
\]
Interpretation: ranking quality and evidence coverage are both incomplete, but in different ways.

\section*{Worked Example (Failure Localization)}
Pipeline A improves final answer score from 68 to 72.  
But decomposition reveals:
\begin{itemize}[leftmargin=1.5em]
  \item retrieval recall@5 dropped from 0.81 to 0.70,
  \item generation conditioned on gold context improved substantially.
\end{itemize}
Interpretation: generator improved, retriever regressed.  
Action: fix retrieval stage rather than over-tuning generation prompts.

\section*{Intuition Checkpoints}
\begin{itemize}[leftmargin=1.5em]
  \item Better answer fluency can hide retrieval failures.
  \item High retrieval recall can coexist with poor final answers if synthesis is weak.
  \item Metric definitions require explicit edge-case conventions.
\end{itemize}

\section*{Further Refresh}
\begin{itemize}[leftmargin=1.5em]
  \item Embedding geometry: see Chapter 4 refresher.
  \item Calibration/uncertainty interpretation: see Chapter 1 refresher.
  \item Systems throughput constraints: see Chapter 7 refresher.
  \item \textbf{Manning et al., \emph{IIR}} --- BM25, ranking metrics, IR evaluation.
  \item \textbf{FAISS/HNSW papers} --- ANN retrieval behavior and tradeoffs.
\end{itemize}

\section{Week 10 Refresher: Risk Arithmetic and Assurance Logic}

\subsection*{Setting the Scene}
Governance math is not abstract compliance language; it is decision logic under uncertainty. The field evolved from safety engineering (hazards, layered controls, assurance cases) into AI governance artifacts. The critical insight is traceability: every risk claim must be connected to measurable evidence and triggerable controls.

\subsection*{Notation Introduced in This Chapter}
\begin{itemize}[leftmargin=1.5em]
  \item $R$: baseline risk score.
  \item $R_{\text{res}}$: residual risk after controls.
  \item $\eta_i$: efficacy of control $i$.
  \item $m_t$: monitored metric at time $t$.
  \item $m_0$: baseline metric value.
  \item $\epsilon$: escalation threshold.
\end{itemize}

\subsection*{What You Need To Recall Precisely}
\begin{itemize}[leftmargin=1.5em]
  \item \textbf{Risk decomposition:} severity and likelihood are different axes with different uncertainty.
  \item \textbf{Layered controls:} multiplicative residual-risk models assume forms of independence.
  \item \textbf{Trigger logic:} indicator functions implement threshold decisions.
  \item \textbf{Assurance structure:} claim $\to$ evidence $\to$ control action $\to$ artifact.
  \item \textbf{Correlation caveat:} multiplicative risk-reduction models are optimistic under shared-cause failures.
\end{itemize}

\subsection*{How These Results Build on Each Other}
\begin{enumerate}[leftmargin=1.5em]
  \item Define risk factors and their measurement status.
  \item Model baseline and post-control residual risk.
  \item Apply threshold policy to operational metrics.
  \item Attach each decision to auditable evidence.
\end{enumerate}

\subsection*{Reminder Glossary (Term by Term)}
\begin{itemize}[leftmargin=1.5em]
  \item \textbf{Hazard:} undesirable outcome class to prevent/mitigate.
  \item \textbf{Control efficacy:} estimated fractional risk reduction from a control.
  \item \textbf{Residual risk:} remaining risk after controls.
  \item \textbf{Threshold policy:} decision rule on monitored metrics.
  \item \textbf{Assurance artifact:} evidence item supporting a claim.
\end{itemize}

\subsection*{Worked Mini Example (Single vs Two Controls)}
If baseline risk score is $R=1.0$, and control efficacies are $\eta_1=0.4$, $\eta_2=0.3$:
\[
R_1=R(1-\eta_1)=0.6,\qquad
R_{12}=R(1-\eta_1)(1-\eta_2)=0.42.
\]
Interpretation: layered controls reduce residual risk multiplicatively under the model assumptions.

\subsection*{Worked Example (Threshold and Escalation)}
Suppose drift metric baseline is $m_0=0.12$ and policy threshold is $\epsilon=0.05$.  
Current observed $m_t=0.19$ gives deviation $|m_t-m_0|=0.07>\epsilon$.  
Escalation rule triggers high-review mode.  
Reasoning chain: measured metric -> threshold test -> policy action -> logged evidence.

\subsection*{Intuition Checkpoints}
\begin{itemize}[leftmargin=1.5em]
  \item Multiplicative control models can be optimistic if failures are correlated.
  \item Governance thresholds are policy design decisions, not immutable laws.
  \item Audit readiness is about traceable evidence chains, not only numeric scores.
\end{itemize}

\subsection*{Further Refresh}
\begin{itemize}[leftmargin=1.5em]
  \item Metric reliability and uncertainty: see Week 1 refresher.
  \item Evaluation decomposition and harness logic: see Week 9 refresher.
  \item \textbf{Leveson, \emph{Engineering a Safer World}} --- hazard/control argument structure.
  \item \textbf{NIST AI RMF} --- operational governance controls and documentation patterns.
\end{itemize}

\chapter{Chapter 12 Refresher: Modern LLM Access --- APIs, SDKs, Agents, and the Model Context Protocol --- Latency, Queues, and Reliability Budgets}

\section*{Setting the Scene}
APIs expose stochastic models inside stochastic infrastructure: arrivals, service times, retries, and pricing multiply Preface Steps~5--7 expectations into tail experiences users actually feel.

\paragraph{Step A --- Utilisation as latent probability.}
\textbf{Problem.} Saturated GPUs inflate latency although average load looks fine.
\textbf{Insight.} \(\rho=\lambda/\mu\) nearing 1 blows variance (Preface Step~6)---queue waits explode superlinearly.
\textbf{Lineage.} Kendall queue notation; Kleinrock textbooks.
\textbf{Mental model.} Treat \(\rho\) like criticality in statistical mechanics---small deltas matter near phase transitions.

\paragraph{Step B --- Retries compound load.}
\textbf{Problem.} Clients retry on flakes.
\textbf{Insight.} Geometric retry counts lift effective \(\lambda\)---same algebra as Preface Step~4 inverted conditioning stories but operational.

\paragraph{Step C --- Guards cap rare disasters.}
\textbf{Problem.} Agents might loop unbounded.
\textbf{Insight.} Hard budgets reinstate finite \(\Omega\) partitions compatible with Step~1 measurability mindset.

Probability-minded reliability engineers read LM dashboards the same way traders read risk tapes---moments plus tails.

\section*{Notation Introduced in This Chapter}
\begin{itemize}[leftmargin=1.5em]
  \item $\lambda$: request arrival rate.
  \item $\mu$: service rate.
  \item $\rho=\lambda/\mu$: utilization.
  \item $\mathbb{E}[W]$: expected waiting time.
  \item $p_f$: per-attempt failure probability.
  \item $\mathbb{E}[K]$: expected number of attempts.
\end{itemize}

\section*{What You Need To Recall Precisely}
\begin{itemize}[leftmargin=1.5em]
  \item \textbf{Latency decomposition:} prefill latency plus per-output-token decode latency.
  \item \textbf{Queueing baseline (M/M/1).}
  Assume Poisson arrivals rate \(\lambda\), exponential service rate \(\mu\) per busy server, single server, infinite buffer (steady-state regime \(\lambda<\mu\)). Expected waiting time in queue satisfies
  \[
    \mathbb{E}[W]=\frac{1}{\mu-\lambda},
  \]
  blowing up as \(\rho=\lambda/\mu\uparrow 1$: tiny overload deltas dominate latency tails---engineering intuition beats pretending arrivals are perfectly deterministic.

  \item \textbf{Retry expectation.}
  If each attempt fails independently with probability \(p_f\) and retries repeat until first success (no cap), successes follow geometric distribution with success probability \(1-p_f\). Expected attempts \(\mathbb{E}[K]=\sum_{n\ge1} n\, p_f^{n-1}(1-p_f)=1/(1-p_f)\). Even modest \(p_f\) inflates load multipliers \(\mathbb{E}[K]\) nonlinearly.
  \item \textbf{Token cost model:} base cost is linear in tokens, expected cost depends on retries and aborts.
  \item \textbf{Termination guards:} bounded steps/tokens/cost enforce finite agent behavior.
\end{itemize}

\section*{How These Results Build on Each Other}
\begin{enumerate}[leftmargin=1.5em]
  \item Estimate traffic intensity $(\lambda/\mu)$.
  \item Predict queue delay and saturation risk.
  \item Incorporate retries into effective load and expected cost.
  \item Enforce runtime guards for bounded execution.
\end{enumerate}

\section*{Reminder Glossary (Term by Term)}
\begin{itemize}[leftmargin=1.5em]
  \item \textbf{Arrival rate $\lambda$:} requests per unit time.
  \item \textbf{Service rate $\mu$:} processing capacity per unit time.
  \item \textbf{Utilization $\rho$:} $\lambda/\mu$.
  \item \textbf{Retry multiplier:} expected attempts per successful request.
  \item \textbf{Budget guard:} hard bound on cost/tokens/steps.
\end{itemize}

\section*{Worked Mini Example (Retries)}
Let per-attempt failure probability be $p_f=0.2$ and assume independent retries with no cap.
\[
\mathbb{E}[K]=\frac{1}{1-p_f}=\frac{1}{0.8}=1.25.
\]
If average successful request cost is \$0.04, expected cost with retries is roughly
\[
0.04\times 1.25 = \$0.05.
\]

\section*{Worked Example (Queue and Retry Coupling)}
If base arrival rate is 80 req/s and retries add 25\% expected attempts, effective arrival becomes:
\[
\lambda_{\text{eff}} = 80 \times 1.25 = 100\ \text{req/s}.
\]
If service rate is $\mu=110$ req/s, utilization rises from $0.73$ to $0.91$.  
Interpretation: retry policy can push system into tail-latency regime even when baseline looked safe.

\section*{Intuition Checkpoints}
\begin{itemize}[leftmargin=1.5em]
  \item Systems can meet average SLOs while failing badly in tails.
  \item Queueing approximations are useful planning tools, not exact traffic simulators.
  \item Reliability controls (backoff, circuit breaking) alter both latency and cost.
\end{itemize}

\section*{Further Refresh}
\begin{itemize}[leftmargin=1.5em]
  \item Systems scaling and bottlenecks: see Chapter 7 refresher.
  \item Evaluation and trust boundaries: see Chapter 10 refresher.
  \item Risk thresholding and controls: see Chapter 11 refresher.
  \item \textbf{Harchol-Balter} --- queueing and latency scaling.
  \item \textbf{SRE workbooks} --- retries/backoff/overload control design.
\end{itemize}

\section{Chapter 13 Refresher: A Human-Centric Model for Understanding LLM Errors --- Diagnostic Decomposition and Multi-Objective Reasoning}

\subsection*{Setting the Scene}
The closing chapter consolidates mathematical judgment rather than introducing a new heavy formalism. The field moved from single-score evaluation toward structured diagnosis because one scalar cannot describe heterogeneous failure causes. The critical insight is decomposition before intervention.

\subsection*{Notation Introduced in This Chapter}
\begin{itemize}[leftmargin=1.5em]
  \item Axis labels: structure, meaning, context, grounding, environment.
  \item Conditional slice: subset of cases defined by context/task conditions.
  \item Tradeoff surface: set of feasible points across accuracy, safety, latency, and cost.
\end{itemize}

\subsection*{What You Need To Recall Precisely}
\begin{itemize}[leftmargin=1.5em]
  \item \textbf{Failure-axis decomposition:} structure, meaning, context, grounding, environment are separate tests.
  \item \textbf{Conditional error analysis:} inspect slices and contexts, not only aggregate rates.
  \item \textbf{Correlation vs intervention:} observed association does not imply remediation leverage.
  \item \textbf{Multi-objective tradeoffs:} optimize over accuracy, safety, latency, and cost jointly.
  \item \textbf{Diagnostic ordering:} simple structural failures should be ruled out before attributing deep semantic or governance causes.
\end{itemize}

\subsection*{How These Results Build on Each Other}
\begin{enumerate}[leftmargin=1.5em]
  \item Observe failure instance and classify candidate axes.
  \item Check low-level structural and context conditions first.
  \item Test grounding and environment causes next.
  \item Select intervention matched to diagnosed failure pathway.
\end{enumerate}

\subsection*{Reminder Glossary (Term by Term)}
\begin{itemize}[leftmargin=1.5em]
  \item \textbf{Structure failure:} format/syntax/protocol violation.
  \item \textbf{Meaning failure:} semantic mismatch despite fluent output.
  \item \textbf{Context failure:} missing or misused prompt/task constraints.
  \item \textbf{Grounding failure:} unsupported factual claims.
  \item \textbf{Environment failure:} external system/tool mismatch.
\end{itemize}

\subsection*{Worked Mini Example (Axis Decision)}
Suppose an answer is fluent but cites a nonexistent document.
\begin{itemize}[leftmargin=1.5em]
  \item \textbf{Structure:} pass (syntax and format are valid).
  \item \textbf{Meaning:} partial pass (language coherent).
  \item \textbf{Grounding:} fail (claim cannot be verified against source evidence).
\end{itemize}
The corrective action should prioritize grounding controls (retrieval/citation verification), not merely prompt rephrasing.

\subsection*{Worked Example (Axis-Matched Intervention)}
Case: answer includes correct format and plausible language but wrong regulation citation.
\begin{itemize}[leftmargin=1.5em]
  \item Structure: pass.
  \item Meaning: partial pass.
  \item Grounding: fail (citation not verifiable).
\end{itemize}
Intervention path:
\begin{enumerate}[leftmargin=1.5em]
  \item add retrieval requirement from authoritative corpus,
  \item enforce citation-span verification before final response,
  \item reject answer if citation cannot be matched.
\end{enumerate}
Reasoning: this is not primarily a prompt-style failure; it is evidence-path failure.

\subsection*{Intuition Checkpoints}
\begin{itemize}[leftmargin=1.5em]
  \item Calibration and uncertainty: see Chapter 1 refresher.
  \item A plausible answer can still be ungrounded.
  \item One failure can belong to multiple axes; diagnosis is often multi-label.
  \item The strongest fix is the one that changes the relevant information pathway.
\end{itemize}

\subsection*{Further Refresh}
\begin{itemize}[leftmargin=1.5em]
  \item Evaluation slicing and retrieval grounding: see Chapter 10 refresher.
  \item Risk and assurance framing: see Chapter 11 refresher.
  \item Runtime constraints and failure budgets: see Chapter 12 refresher.
  \item \textbf{Murphy} --- probabilistic uncertainty reasoning.
  \item \textbf{Error analysis/evaluation notes} --- slicing and diagnostic methodology.
\end{itemize}


\backmatter

\end{document}
